Este informe aborda la planificación diaria de la operación de los buses eléctricos de la Red Metropolitana de Movilidad de Santiago, un problema de asignación de vehículos eléctricos conocido como E-VSP. Para un día laboral con 64.502 viajes programados, se decide qué bus realiza cada viaje y cuándo y dónde se recarga, con el objetivo de minimizar el costo total de operación, considerando la autonomía limitada de los buses y la disponibilidad de puestos de carga en los electroterminales. Dada la escala del problema, se propone una metodología en cuatro etapas, que asigna las rutas a electroterminales, construye las jornadas de cada bus permitiendo que atienda viajes de rutas distintas, corta y reencadena esas jornadas según la batería, y programa las recargas mediante un modelo de optimización entera mixta.

Se construyó un caso base para toda la red y se evaluó la propuesta completa en el electroterminal La Reina. En toda la red, permitir que un mismo bus atienda viajes de rutas distintas (\textit{interlining}) reduce el costo del caso base en 5,2\%. Además, incorporar la batería aumenta la flota entre 53\% y 70\% respecto de las jornadas que la ignoran. En La Reina, la propuesta completa reduce el costo diario en 28,6\% y la flota en 32,4\% frente al caso base con carga reactiva, y el ahorro proviene principalmente de reutilizar buses y de disminuir los buses de reserva, más que de cargar en horas más baratas. Estos resultados son preliminares, pues el programa de carga supera brevemente los puestos disponibles, persisten cadenas de buses que exceden las 24 h y la propuesta aún no se evalúa en toda la red. Por ende, el trabajo futuro se centra en extender la propuesta a toda la red y corregir estas limitaciones.
